\subsection{Формула Тейлора, остаточный член в форме Лагранжа.}

\subsubsection*{Построение многочлена (Определение коэффициентов)}

Пусть $P_n(x)$ — многочлен степени $n$:
\[ P_n(x) = p_0 + p_1 x + p_2 x^2 + \dots + p_n x^n \]

Найдем его коэффициенты $p_k$, выразив их через значения производных в нуле.
Дифференцируем многочлен:
\begin{align*}
    P_n(0) &= p_0 \\
    P'_n(x) &= 1 \cdot p_1 + 2p_2 x + 3p_3 x^2 + \dots + n p_n x^{n-1} \\
    P'_n(0) &= 1 \cdot p_1 \implies p_1 = \frac{P'_n(0)}{1!} \\
    P''_n(0) &= 2 \cdot 1 \cdot p_2 \implies p_2 = \frac{P''_n(0)}{2!} \\
    &\dots \\
    P^{(k)}_n(0) &= k! \, p_k \implies p_k = \frac{P^{(k)}_n(0)}{k!}
\end{align*}

Таким образом, многочлен можно записать как:
\[ P_n(x) = P_n(0) + \frac{P'_n(0)}{1!}x + \dots + \frac{P^{(n)}_n(0)}{n!}x^n \]

\subsubsection*{Формула Тейлора для произвольной точки $a$}

Разложим многочлен по степеням $(x-a)$. Сделаем замену переменной $x = (x-a) + a$.
\[ P_n(x) = \tilde{p}_0 + \tilde{p}_1(x-a) + \dots + \tilde{p}_n(x-a)^n \]
Аналогично вычисляя производные в точке $x=a$, получаем:
\[ \tilde{p}_k = \frac{P^{(k)}_n(a)}{k!} \]

\begin{tcolorbox}[title=Определение]
    \textbf{Многочленом Тейлора} для функции $f(x)$, имеющей производные до порядка $n$ в точке $a$, называется многочлен:
    \[ P_n(x, a) = f(a) + \frac{f'(a)}{1!}(x-a) + \dots + \frac{f^{(n)}(a)}{n!}(x-a)^n = \sum_{k=0}^n \frac{f^{(k)}(a)}{k!}(x-a)^k \]
\end{tcolorbox}
Мы пишем $f(x) \approx P_n(x, a)$ (или ставим знак соответствия).

\subsubsection*{Остаточный член}

Чтобы поставить знак равенства, введем разность между функцией и её многочленом Тейлора.
\[ f(x) = P_n(x, a) + r_n(x, a) \]
Где $r_n(x, a)$ — \textbf{остаточный член}.

\textbf{Свойства остаточного члена в точке $a$:}
Так как $P_n(x)$ построен так, что его производные в точке $a$ совпадают с производными $f(x)$, то:
\[ r_n(a) = f(a) - P_n(a) = 0 \]
\[ r'_n(a) = f'(a) - P'_n(a) = 0 \]
\[ \dots \]
\[ r^{(k)}_n(a) = 0 \quad \text{для всех } k = 0, \dots, n. \]

\subsubsection*{Теорема об остаточном члене (Форма Лагранжа)}

Используем вспомогательные функции (для сведения к обобщенной теореме Коши из):
\begin{itemize}
    \item $\varphi(x) := r_{n-1}(x, a)$ (остаток для многочлена степени $n-1$)
    \item $\psi(x) := (x-a)^n$
\end{itemize}
Обе функции и их производные до $(n-1)$-го порядка равны 0 в точке $a$.

\begin{tcolorbox}[title=Теорема (Остаток в форме Лагранжа)]
    Пусть $f(x) \in C[a, b]$ и имеет производные до порядка $n$ на $(a, b)$ (или $(n-1)$ в замкнутой области, согласно записям).
    
    Для любого $x \in [a, b]$ существует точка $\xi \in (a, x)$, такая что справедливо разложение:
    \[ f(x) = \sum_{k=0}^{n-1} \frac{f^{(k)}(a)}{k!}(x-a)^k + \frac{f^{(n)}(\xi)}{n!}(x-a)^n \]
\end{tcolorbox}

Это \textbf{глобальная форма} формулы Тейлора.
Остаточный член:
\[ r_{n-1}(x) = \frac{f^{(n)}(\xi)}{n!}(x-a)^n \]


